\chapter*{Motivație}
\addcontentsline{toc}{chapter}{Motivație}

Ideea acestei aplicații a apărut natural, din observațiile pe care le-am făcut
în anii de studiu la Facultatea de Automatică și Calculatoare din Iași. Ca student,
am interacționat constant cu sistemele informatice ale universității și am realizat că,
deși digitalizarea este prezentă, există încă un loc de îmbunătățire între platformele
oficiale de curs și nevoile practice de la laborator.

În prezent, platforma Moodle este standardul pentru gestionarea cursurilor și a temelor.
Este un sistem robust, însă, fiind o soluție generalistă, nu acoperă în totalitate dinamica specifică
a activităților de laborator. Am observat o fragmentare a datelor pe parcursul fiecărui semestru,
informații precum prezențele, punctajele intermediare sau organizarea grupelor ajungând să fie dispersate
în fișiere Excel locale, mesaje pe email sau note pe hârtie. Această metodă de lucru în paralel 
generează o lipsă de transparență care îngreunează procesul educațional.

Motivația principală a proiectului Academo a fost eliminarea acestor bariere prin crearea
unui punct unic de acces. Pentru student, cea mai importantă este certitudinea informației.
Prin Academo, mi-am propus să ofer un loc centralizat unde parcursul academic poate
fi verificat instantaneu: prezențele, activitățile bifate și observațiile relevante
sunt vizibile imediat. Această transparență permite studentului să se concentreze pe
învățare, fără a mai avea grija monitorizării manuale a propriului progres. În plus,
am dorit ca aplicația să servească și ca un hub logistic, oferind posibilitatea de a
descărca orarul și programarea examenelor direct din interfață. Astfel, utilizatorii nu mai 
sunt nevoiți să acceseze separat site-ul universității, având toate informațiile esențiale 
într-un singur loc.

În același timp, aplicația este gândită să simplifice munca profesorilor prin oferirea unui 
control total asupra activității de la grupă. Un aspect esențial este posibilitatea de a 
actualiza în timp real conținutul laboratoarelor, permițând adaptarea sarcinilor de lucru. 
De asemenea, centralizarea listelor de studenți în tabele intuitive și accesul la 
un istoric complet al activităților oferă profesorului o viziune clară asupra evoluției semestrului. 
Posibilitatea de a exporta aceste situații centralizate în format Excel transformă raportarea finală 
dintr-o sarcină manuală obositoare într-o acțiune de câteva secunde.

Din punct de vedere tehnic, realizarea acestei platforme a reprezentat oportunitatea de a-mi consolida 
cunoștințele de inginerie software. Am ales să utilizez tehnologii moderne precum Angular și PHP pentru 
a construi un sistem performant și ușor de utilizat. În final, Academo este rezultatul dorinței de a 
aduce mai multă rigoare în mediul universitar, demonstrând cum tehnologia poate simplifica interacțiunea 
dintre studenți și profesori prin centralizare și transparență.

\newpage

\chapter*{Introducere}
\addcontentsline{toc}{chapter}{Introducere}

Digitalizarea în mediul universitar este un proces esențial, dar care nu este întotdeauna complet. 
În facultatea noastră, deși platformele oficiale acoperă bine partea de cursuri, activitățile 
practice de laborator rămân adesea într-o zonă gri. Informațiile despre prezențe, organizarea 
grupelor sau punctajele intermediare ajung să fie dispersate în fișiere Excel salvate local sau 
mesaje pe email. Această fragmentare a datelor face ca procesul administrativ să fie greoi și ineficient.

Această situație creează o lipsă de transparență care afectează ambele părți implicate. Din perspectiva 
studentului, este dificil să ai o imagine actualizată asupra propriului parcurs fără a apela direct 
la profesor. De cealaltă parte, profesorii consumă timp prețios sincronizând manual date din surse 
diferite, ceea ce încetinește procesul didactic și crește riscul de a introduce greșit informații 
importante în sistemul final.

Lucrarea de față propune o soluție prin dezvoltarea aplicației Academo, o platformă web gândită să 
funcționeze ca un punct central de gestiune. Scopul proiectului este să ofere un instrument complementar, 
rapid și intuitiv, adaptat pentru nevoile laboratoarelor. Prin centralizarea fluxurilor într-o singură 
bază de date, aplicația asigură o sursă unică de adevăr pentru toți utilizatorii, oferind tabele clare 
pentru managementul studenților și al grupelor.

Din punct de vedere funcțional, aplicația permite profesorilor să gestioneze dinamic desfășurarea orelor. 
Aceștia pot actualiza conținutul laboratoarelor în timp real și pot consulta un istoric detaliat al 
activităților realizate, având astfel o evidență clară a tuturor modificărilor. Interfața oferă 
instrumente pentru marcarea rapidă a prezenței și exportul statisticilor 
în format Excel, automatizând sarcinile repetitive. De asemenea, atât studenții, cât și profesorii 
pot descărca orarul și programarea examenelor direct din aplicație, eliminând necesitatea de a accesa 
site-ul universității pentru aceste detalii logistice de rutină.

Din punct de vedere tehnic, am ales un stack tehnologic modern. Pentru interfață am utilizat Angular 19, 
profitând de noile funcționalități precum Signals pentru o reactivitate rapidă. Pentru partea de server 
am utilizat limbajul PHP, implementând un API RESTful care gestionează logica de business. Arhitectura 
este una decuplată, unde frontend-ul și backend-ul comunică prin HTTP folosind mesaje JSON, 
ceea ce face ca sistemul să fie ușor de întreținut.

În final, Academo reprezintă o soluție tehnică la o problemă reală. Prin oferirea accesului în timp real 
la date și centralizarea tuturor instrumentelor de lucru și a informațiilor administrative într-un singur 
loc, proiectul își propune să elimine barierele birocratice pentru ca studenții și profesorii să se poată 
concentra pe procesul de învățare.

\newpage